\subsection{Sprint 4}

A Sprint 4 começou com uma das tarefas mais estruturais que assumi no projeto: a substituição completa dos dados mockados do fluxo de autenticação por uma integração real com o backend, abrangendo desde a persistência segura do \textit{token} de autenticação até o ajuste dos \textit{schemas} de \textit{Manager}, \textit{Salesman}, \textit{Company}, \textit{Meeting} e \textit{Dashboard} para refletirem com fidelidade os contratos reais da API. Essa tarefa, embora tecnicamente trabalhosa, era um pré-requisito indispensável para que o restante do time pudesse desenvolver sobre dados confiáveis em vez de respostas simuladas.

Em seguida, voltei a exercer meu papel de revisão arquitetural validando e integrando dois \textit{pull requests} relevantes para a saúde do projeto a longo prazo: a introdução de um sistema de notificações \textit{toast} compartilhado entre as páginas e uma refatoração ampla de arquivos que estavam fora da estrutura de pastas definida na arquitetura do frontend. Também contribuí de forma direta com a qualidade do projeto ao criar, do zero, uma suíte de testes \textit{end-to-end} com Playwright \cite{Playwright}, cobrindo os fluxos completos de login, \textit{dashboard}, empresas, gestores, vendedores e reuniões (incluindo todas as abas de detalhe) contra um \textit{backend} real com dados \textit{seedados} via Docker, além de ajustar a configuração do Vitest \cite{Vitest} para que essas especificações não fossem executadas indevidamente pelos testes unitários.

Outro problema enfrentado na sprint, dessa vez de origem externa ao frontend, foi a constatação de que cinco \textit{endpoints} do \textit{backend} necessários para a integração das telas ainda estavam faltando muito próximo da entrega, pois não haviam sido testados a tempo durante o final de semana que antecedia o prazo. Como poucos colegas estavam disponíveis para resolver a pendência naquele momento, a situação terminou sendo resolvida por um colega de outra turma que assumiu a implementação dos \textit{endpoints} faltantes. Esse episódio ficou como lição para minhas próprias entregas em \ac{ages} futuras: a importância de testar o projeto de ponta a ponta com antecedência e não apenas próximo ao prazo de entrega, já que problemas de última hora encontrados sem tempo ou disponibilidade de equipe para serem resolvidos colocam em risco a entrega como um todo.

O momento mais desafiador da sprint ocorreu próximo à sua entrega, quando uma quantidade expressiva de funcionalidades desenvolvidas em paralelo por diferentes colegas, como a edição de vendedores pelo gestor, o acesso aos detalhes de vendedores e o aprimoramento do painel e das telas de reuniões do gestor e do vendedor, convergiu para a \textit{branch} de desenvolvimento em um curto intervalo de tempo. Essa integração concentrada revelou uma série de inconsistências que não haviam sido percebidas isoladamente em cada \textit{pull request}: campos de empresa que deveriam ser removidos e não apenas desabilitados durante a edição, buscas que não ignoravam acentuação, modais de cadastro vulneráveis a duplo envio, métricas de sentimento ainda calculadas a partir de valores fixos em vez de dados reais, e uma lacuna no filtro de reuniões no lado do cliente. Diante disso, assumi a responsabilidade de mapear e corrigir essas inconsistências em uma sequência intensa de correções, unificando a fonte de dados de vendedores ativos, conectando os campos de endereço e telefone da empresa aos dados reais e documentando as escalas de sentimento utilizadas, sempre acompanhando essas correções da atualização dos testes correspondentes.

Essa experiência deixou uma lição clara sobre os limites da revisão de \textit{pull requests} isolados: problemas de integração entre funcionalidades desenvolvidas em paralelo só se tornam visíveis quando o sistema é observado como um todo, o que confirma o valor da suíte de testes \textit{end-to-end} construída no início da sprint e aponta para a necessidade de integrações mais frequentes e incrementais nas próximas sprints, em vez de grandes convergências de última hora. Nesse sentido, o ponto levantado na retrospectiva da Sprint 3 sobre a necessidade de fechar as \textit{tasks} com mais antecedência foi parcialmente endereçado pelos colegas de \ac{ages} IV, que passaram a estabelecer prazos mais claros para a conclusão das tarefas. Isso não eliminou por completo a abertura de \textit{pull requests} na última semana da sprint, o que é até certo ponto natural dado o ritmo do desenvolvimento, mas fez com que a maioria das \textit{tasks} com datas definidas fosse concluída dentro do prazo, permitindo que boa parte dos \textit{pull requests} fosse revisada com calma ao longo da sprint em vez de se concentrar inteiramente em seus últimos dias. Essa melhoria, ainda que incompleta, foi reconhecida pelo time na retrospectiva de encerramento da Sprint 4. Com as inconsistências do painel do gestor e do vendedor sanadas e a base de autenticação e notificações consolidada, o time se encontra em uma posição mais sólida para finalizar as funcionalidades remanescentes do produto e preparar as entregas finais do projeto SalesPilot.
